\SourcePage{57}
\SGAsection{3}{关于 K(X) 的函子性的一般讨论}
若 $X$ 是\Term{algebraic-space}{代数空间}，为简便记 $K(X)=K(\mathscr T(X))$；当 $X$ \Term{quasi-projective}{拟射影}且\Term{nonsingular}{非奇异}时，依\Term{theorem}{定理} 2.3 的推论 2，把此群与 $K(\mathscr S(X))$ 等同。
\SourcePage{58}
由于 $K(\mathscr S(X))$ 不仅是\Term{abelian-group}{阿贝尔群}，而且是 \Term{lambda-ring}{$\lambda$-环}，传递结构便使 $K(X)$ 也成为 \Term{lambda-ring}{$\lambda$-环}，因而带有\Term{multiplication}{乘法}结构及映射 $\lambda^i$。必须不经过\Term{vector-bundle}{向量丛}，直接定义这些结构；下面就来这样做。

若 $E$ 是 $X$ 上的\Term{vector-bundle}{向量丛}，以 $\gamma(E)$ \Term{representation}{表示}它在 $K(X)$ 中的类；若 $F$ 是 $X$ 上的\Term{coherent-algebraic-sheaf}{凝聚代数层}，同样以 $\gamma(F)$ \Term{representation}{表示}它在 $K(X)$ 中的类。若 $Y$ 是 $X$ 的\Term{irreducible}{不可约}\Term{subvariety}{子簇}，置 $\gamma(Y)=\gamma(\mathcal O_Y)$，其中 $\mathcal O_Y$ 是 $Y$ 的\Term{local-ring}{局部环}\Term{sheaf}{层}，视为 $X$ 上的\Term{coherent-algebraic-sheaf}{凝聚代数层}。由\Term{linear}{线性}性把\Term{signature}{符号} $\gamma(Y)$ 定义到 $X$ 上任意\Term{algebraic-cycle}{代数闭链} $Y$。可能混淆时，写 $\gamma_X$ 代替 $\gamma$。

设 $X$ 为任意\Term{algebraic-space}{代数空间}。在 $K(X)$ 上引入\Term{filtration}{滤过}：以 $K(X)^i$ \Term{representation}{表示}由所有满足 $\dim\operatorname{supp}F\leq n-i$ 的 $\gamma(F)$ \Term{generated-subgroup}{生成的子群}，这里 $n$ 是 $X$ 的\Term{dimension}{维数}。由 Serre 的“\Term{devissage-lemma}{逐层分解引理}”（采用 [3] 所给的\Term{form}{形式}），得到如下结果。
\begin{SGAresult}{命题}{2.6}
群 $K(X)^i$ 由 $\gamma(Y)$ 生成，其中 $Y$ 遍历 $X$ 的\Term{dimension}{维数} $\leq n-i$ 的\Term{subvariety}{子簇}。若 $X$ \Term{irreducible}{不可约}，则 $G^0(K(X))=\mathbb Z$。
\end{SGAresult}
最后这个\Term{isomorphism}{同构}由\Term{coherent-algebraic-sheaf}{凝聚代数层}的\Term{rank}{秩}所定义的\Term{homomorphism}{同态} $K(X)\to\mathbb Z$ 给出。

现在设 $X$ \Term{nonsingular}{非奇异}。仿照第 2 节开头的讨论，
\SourcePage{59}
在 $K(X)$ 中定义\Term{multiplication}{乘法}如下：
\[
\gamma(F)\gamma(G)=\sum_{i\geq0}(-1)^i\gamma(\operatorname{Tor}_i(F,G)).\tag{2.15}
\]
这里的 $\operatorname{Tor}_i$ 当然是相对于 $X$ 的\Term{local-ring}{局部环}取的。$\operatorname{Tor}_i$ 的\Term{exact-sequence}{正合列}表明，公式 (2.15) 确实定义了从 $K(X)\times K(X)$ 到 $K(X)$ 的\Term{bilinear-map}{双加性映射}；\Term{associative-law}{结合律}由多重 Tor 的\Term{spectral-sequence}{谱序列}得出。这样也可验证，多个因子的乘积可以由下式计算：
\[
\gamma(F_1)\cdots\gamma(F_p)=\sum_{i\geq0}(-1)^i\gamma(\operatorname{Tor}_i(F_1,\ldots,F_p)).\tag{2.16}
\]
右端的 $\operatorname{Tor}_i$ 是“同时 Tor”，通过对各变元 $F_i$ 同时作\Term{projective-resolution}{投射消解}来计算。

若在 (2.15) 中设 $G$ \Term{locally-free}{局部自由}，便得到
\[
\gamma(F)\gamma(G)=\gamma(F\otimes G).\tag{2.15 bis}
\]
这特别表明，$\gamma(\mathcal O_X)$ 是 $K(X)$ 的\Term{identity-element}{单位元}，并且\Term{canonical}{典范}\Term{homomorphism}{同态} $K(\mathscr S(X))\to K(X)$ 是\Term{ring-homomorphism}{环同态}。

公式 (2.15) 显然受到\Term{intersection-theory}{交理论}中 Serre 的“Tor 公式”的启发。该公式证明了下述命题的第二部分；第一部分直接由 (2.15) 和\Term{local-scheme}{局部} Tor 的计算得出。
\begin{SGAresult}{命题}{2.7}
设 $X$ 是\Term{nonsingular-variety}{非奇异簇}，$Y,Z$ 是 $X$ 中两个\Term{algebraic-cycle}{代数闭链}。若 $Y,Z$ 处处\Term{transverse-intersection}{横截相交}，则
\[
\gamma(Y)\gamma(Z)=\gamma(Y\cdot Z).\tag{2.17}
\]
若 $Y,Z$ \Term{homogeneous}{齐次}，\Term{dimension}{维数}分别为 $n-i,n-j$，且它们的交没有\Term{excess-component}{超额分支}，则
\SourcePage{60}
\[
\gamma(Y)\gamma(Z)\equiv\gamma(Y\cdot Z)\pmod {K(X)^{i+j+1}}.\tag{2.18}
\]
\end{SGAresult}
注意，(2.18) 中\Term{generic}{一般}并无等式。第 4 节将说明：若 $X$ \Term{quasi-projective}{拟射影}且\Term{nonsingular}{非奇异}，则 $K(X)$ 的\Term{ring-structure}{环结构}与其\Term{filtration}{滤过}相容，而且\Term{associated-graded-ring}{相伴分次环} $G(K(X))$ 可与 $A(X)$ 模去一个\Term{torsion-subgroup}{挠子群}的商等同（也许该\Term{subgroup}{子群}总是 $0$\footnote{并非如此；反例见本讨论班第 XIV 讲 4.7。}）。

只有在域 $k$ 的\Term{characteristic}{特征}为 $0$ 时，我才知道如何直接用\Term{sheaf}{层}定义 $\lambda^i(F)$（这正是妨碍我们在 $k$ 的\Term{characteristic}{特征}非零时证明 \Term{riemann-roch-theorem}{Riemann--Roch 定理}的原因）。设 $F$ 是 $X$ 上的\Term{coherent-algebraic-sheaf}{凝聚代数层}，$p$ 个字母的\Term{permutation-group}{置换群} $\mathfrak S_p$ 在其上\Term{action}{作用}；以 $(F)^{\mathrm{alt}}$ \Term{representation}{表示}满足对每个 $s\in\mathfrak S_p$ 都有 $sx=\varepsilon(s)x$ 的 $x\in F$ 的全体，$\varepsilon(s)$ 为 $s$ 的\Term{signature}{符号}。这显然是 $F$ 的\Term{coherent}{凝聚}\Term{algebra}{代数}子\Term{sheaf}{层}。于是有：
\begin{SGAresult}{\Term{theorem}{定理}}{2.8}
设 $X$ 是\Term{characteristic}{特征} $0$ 的域 $k$ 上的\Term{nonsingular}{非奇异}\Term{quasi-projective}{拟射影}簇，$F$ 是 $X$ 上的\Term{coherent-algebraic-sheaf}{凝聚代数层}。则
\[
\lambda^p(\gamma(F))=\sum_{i\geq0}(-1)^i\gamma\bigl((\operatorname{Tor}_i(F,F,\ldots,F))^{\mathrm{alt}}\bigr).\tag{2.19}
\]
\end{SGAresult}
\textbf{注意。} 群 $\mathfrak S_p$ 显然\Term{action}{作用}在\Term{sheaf}{层} $\operatorname{Tor}_i(F,\ldots,F)$ 上。

\textbf{证明的思路。} 取 $F$ 的一个由\Term{locally-free}{局部自由}\Term{sheaf}{层} $L_i$ 组成的\Term{finite-resolution}{有限消解} $L$，并使用下面的\Term{generic}{一般}\Term{lemma}{引理}。
\SourcePage{61}
\begin{SGAresult}{\Term{lemma}{引理}}{2.9}
设 $L=(L_i)$ 是\Term{graded}{分次}\Term{locally-free}{局部自由}\Term{coherent-algebraic-sheaf}{凝聚代数层}，只有\Term{finite}{有限}个非零\Term{set-component}{分量}。在 $K(X)$ 中有公式
\[
E\bigl((\mathop{\bigotimes}\nolimits^p L)^{\mathrm{alt}}\bigr)=\lambda^p(E(L)).\tag{2.20}
\]
\end{SGAresult}
此式中，将 $\bigotimes^p L$ 视为带有群 $\mathfrak S_p$ \Term{action}{作用}的\Term{sheaf}{层}；\Term{action}{作用}时须计入\Term{graded}{分次}，故按通常方式引进\Term{signature}{符号}，参见 Cartan--Eilenberg。此外，若 $M$ 是 $X$ 上的\Term{graded}{分次}\Term{coherent-algebraic-sheaf}{凝聚代数层}，置
\[
E(M)=\sum_{i\geq0}(-1)^i\gamma(M_i).
\]
\begin{proof}
证明 (2.20)，可归结为在 $K(G)$ 中证明相应公式，其中 $G$ 为\Term{algebraic-group}{代数群}（愿意的话，可取为若干 $\operatorname{GL}(n,k)$ 的乘积），$L$ 是 $k$ 上\Term{finite-dimensional}{有限维} $G$-模（$K(G)$ 的定义见第一章第 2 节例 2）。在这种情形，这个公式已经在\Term{theorem}{定理} 1.4 中使用，用来说明 (1.33) 右端满足与左端相同的函数方程；验证是初等的。
\end{proof}

设 $f$ 是 $X$ 到\Term{algebraic-space}{代数空间} $Y$ 的\Term{morphism}{态射}。\Term{vector-bundle}{向量丛}的\Term{inverse-image}{逆像}定义了 $\lambda$-\Term{ring-homomorphism}{环同态}
\[
f^!:K(\mathscr S(Y))\longrightarrow K(\mathscr S(X)).
\]
把\Term{vector-bundle}{向量丛}与\Term{locally-free}{局部自由}\Term{algebra}{代数}\Term{sheaf}{层}等同，\Term{vector-bundle}{向量丛}的\Term{inverse-image}{逆像}概念便对应于\Term{algebra}{代数}\Term{sheaf}{层}的\Term{inverse-image}{逆像}概念（[3, 第 8 号] 中已使用）。当 $Y$ \Term{nonsingular}{非奇异}时，更\Term{generic}{一般}地定义\Term{group-homomorphism}{群同态}
\[
f^!:K(Y)\longrightarrow K(X),\tag{2.21}
\]
\SourcePage{62}
所用公式含 Tor 的\Term{alternating-sum}{交错和}，受第 2 节开头讨论的启发，读者可自行写出。于是下图交换：
\[
\begin{tikzcd}
K(\mathscr S(Y))\arrow[r,"f^!"]\arrow[d]&K(\mathscr S(X))\arrow[d]\\
K(Y)\arrow[r,"f^!"]&K(X).
\end{tikzcd}
\]
这特别表明，当 $X,Y$ 都\Term{nonsingular}{非奇异}且\Term{quasi-projective}{拟射影}时，在把 $K(X),K(Y)$ 分别与 $K(\mathscr S(X)),K(\mathscr S(Y))$ 等同之后，$f^!$ 的两个定义一致。

另一个无需 $Y$ \Term{nonsingular}{非奇异}便可定义\Term{homomorphism}{同态} (2.21) 的情形，是\Term{functor}{函子} $F\mapsto f^!(F)$ 从 $\mathscr T(Y)$ 到 $\mathscr T(X)$ \Term{exact}{正合}，即 $X$ 的\Term{local-ring}{局部环}是 $Y$ 的\Term{local-ring}{局部环}上的\Term{flat}{平坦}模。此时只须令
\[
f^!(\gamma_Y(F))=\gamma_X(f^!(F)).
\]
特别地，若 $X$ 是 $Y$ 上\Term{local-triviality}{局部平凡}的\Term{fibred-space}{纤维空间}，便属于这种情形（为看出这一点，可归结为乘积 $X=T\times Y$）。又可验证，在这种情况下，若 $Z$ 是 $Y$ 的\Term{irreducible}{不可约}\Term{subvariety}{子簇}，则 $f^!(\mathcal O_Z)=\mathcal O_{f^{-1}(Z)}$，所以
\[
f^!(\gamma_Y(Z))=\gamma_X(f^{-1}(Z)),\tag{2.22}
\]
条件是 $X$ 在 $Y$ 上“\Term{flat}{平坦}”。没有最后这个条件，公式 (2.22) \Term{generic}{一般}不成立；但有下述与命题 2.7 类似的命题。
\begin{SGAresult}{命题}{2.8}
设 $f:X\to Y$ 是\Term{nonsingular-variety}{非奇异簇}之间的\Term{morphism}{态射}，$Z$ 是 $Y$ 上的\Term{algebraic-cycle}{代数闭链}。若 $Z$ 与 $f$ 处处\Term{transverse}{横截}，则公式 (2.22) 成立；若 $Z$ \Term{homogeneous}{齐次}且\Term{codimension}{余维}为 $i$，并且 $f^{-1}(Z)$ 没有\Term{excess-component}{超额分支}，则
\SourcePage{63}
\[
f^!(\gamma_Y(Z))\equiv\gamma_X(f^{-1}(Z))\pmod {K(X)^{i+1}}.\tag{2.23}
\]
\end{SGAresult}
\begin{proof}与命题 2.7 的证明相同。\end{proof}
最后指出，第 4 节将说明：当 $X,Y$ \Term{quasi-projective}{拟射影}且\Term{nonsingular}{非奇异}时，$f^!$ 与 $K(X),K(Y)$ 的\Term{filtration}{滤过}相容。

再设 $X,Y$ 为任意两个\Term{algebraic-space}{代数空间}，$f:X\to Y$ 是\Term{proper}{固有}\Term{morphism}{态射}。对 $X$ 上任意\Term{coherent-algebraic-sheaf}{凝聚代数层} $F$，\Term{direct-image}{直接像}\Term{sheaf}{层} $f_!(F)$ 是\Term{coherent-algebraic-sheaf}{凝聚代数层}，更\Term{generic}{一般}的\Term{sheaf}{层} $R^qf_!(F)$ 也如此（见 [3]）。回忆按定义，$R^qf_!(F)$ 是 $Y$ 上与\Term{presheaf}{预层} $U\mapsto H^q(f^{-1}(U),F)$ 相伴的\Term{sheaf}{层}。特别有
\[
R^0f_!(F)=f_!(F),\qquad R^qf_!(F)=0\quad(q>\dim X),\tag{2.24}
\]
并且，当\Term{open-subset}{开集} $U$ \Term{affine}{仿射}时，
\[
\Gamma(U,R^qf_!(F))=H^q(f^{-1}(U),F).\tag{2.25}
\]
依第 2 节，由下式定义 $K(X)$ 到 $K(Y)$ 的\Term{homomorphism}{同态}，仍记作 $f_!$：
\[
f_!(\gamma_X(F))=\sum_{q\geq0}(-1)^q\gamma_Y(R^qf_!(F)).\tag{2.26}
\]
这样，$K(X)$ 对于\Term{proper}{固有}\Term{morphism}{态射}成为 $X$ 的\Term{covariant-functor}{协变函子}。关系式 $(fg)_!=f_!g_!$（以及先前忘了指出的 $(fg)^!=g^!f^!$）是第 2 节提及的\Term{transitivity}{传递性}关系的特例；
\SourcePage{64}
这些关系来自\Term{composite-functor}{复合函子}的\Term{spectral-sequence}{谱序列}。

记下如下公式，它在 $Y$ \Term{nonsingular}{非奇异}时成立（若 $Y$ 还\Term{quasi-projective}{拟射影}，取 $y$ 为 $Y$ 上某个\Term{vector-bundle}{向量丛}的类，便特别容易验证）：
\[
f_!(x\cdot f^!(y))=f_!(x)\cdot y\qquad(x\in K(X),\ y\in K(Y)).\tag{2.27}
\]
特别地，若 $Y$ 的点在 $f$ 下的\Term{inverse-image}{逆像}\Term{finite}{有限}且包含在一个\Term{affine-open}{仿射开集}中\footnote{最后这个条件显然由第一个条件\Term{pushout}{推出}（作者在 1957 年尚不知道这一点！）。}，则可知\Term{functor}{函子} $F\mapsto f_!(F)$ \Term{exact}{正合}，且对 $q>0$ 有 $R^qf_!(F)=0$；于是 $f_!(\gamma_X(F))=\gamma_Y(f_!(F))$。特别地，若 $X$ 是 $Y$ 的\Term{algebraic-subspace}{代数子空间}，$i$ 是 $Y$ 到 $X$ 的\Term{canonical}{典范}\Term{injective-map}{单射}，则
\[
i_!(\gamma_X(F))=\gamma_Y(F),\tag{2.28}
\]
其中右端把 $F$ 视为 $Y$ 上在 $X$ 外为零的\Term{sheaf}{层}。

\SGAsection{4}{若干技术性结果}
设 $X$ 是\Term{algebraic-space}{代数空间}；第 3 节已见，\Term{projection}{投影} $p:X\times k\to X$ 定义\Term{homomorphism}{同态}
\[
p^!:K(X)\longrightarrow K(X\times k).\tag{2.29}
\]
\begin{SGAresult}{命题}{2.9}\Term{homomorphism}{同态} (2.29) 是\Term{surjective-map}{满射}。\end{SGAresult}
\begin{proof}
对 $n=\dim X$ 归纳；$n<0$ 时断言平凡。设 $n\geq0$；依命题 2.6，只须验证，对于 $X\times k$ 的每个\Term{irreducible}{不可约}\Term{closed-subset}{闭子集} $Y$，$\gamma(Y)$ 都在 $p^!$ 的像内。
\SourcePage{65}
继续论证时只须设 $X$ \Term{irreducible}{不可约}。若 $p(Y)$ 在 $X$ 中不\Term{dense}{稠密}，则断言由归纳假设得到；否则，或者 $Y=X\times k$，断言平凡，或者 $Y$ 是 $X\times k$ 中的\Term{hypersurface}{超曲面}。取由 $X$ 的\Term{nonsingular}{非奇异}点组成的\Term{affine-open}{仿射开集} $U$；若 $D$ 是 $Y$ 在 $U\times k$ 上诱导的\Term{divisor}{除子}，则由\Term{commutative-algebra}{交换代数}知道，它\Term{linear-equivalence}{线性等价}于形如 $D'\times k$ 的\Term{divisor}{除子}，其中 $D'$ 是 $U$ 上的\Term{divisor}{除子}。于是
\[
\gamma_{U\times k}(D)-\gamma_{U\times k}(D'\times k)\in K(U\times k)^2
\]
（命题 2.10 的推论），并且由已经证明的部分，此元素形如 $p^!(x')$，其中 $x'\in K(U)$。所以 $\gamma_{U\times k}(D)=p^!(y')$，其中 $y'=x'+\gamma_U(D')$。但 $y'$ 是某个 $y\in K(X)$ 的\Term{restriction}{限制}（命题 2.10），故 $\gamma_{X\times k}(D)-p^!(y)$ 在 $U\times k$ 上的\Term{restriction}{限制}为零。由命题 2.10，$\gamma_{X\times k}(D)-p^!(y)$ 来自
\[
K(X\times k-U\times k)=K(X'\times k),\qquad X'=X-U
\]
中的某个元素。再用一次归纳假设即完成证明。
\end{proof}
\begin{SGAresult}{推论}{}
设 $X$ 是\Term{algebraic-space}{代数空间}，$p$ 是 $X\times k^n$ 到 $X$ 的\Term{projection}{投影}。则 $p^!:K(X)\to K(X\times k^n)$ 是\Term{surjective-map}{满射}；若 $X$ \Term{nonsingular}{非奇异}，它还是\Term{bijective-map}{双射}。
\end{SGAresult}
\begin{proof}
第一断言对 $n$ 归纳。对于第二断言，令 $f(x)=(x,0)$，故 $pf=1$，$f^!p^!=1$；这证明 $p^!$ 为\Term{injective-map}{单射}。
\end{proof}
\begin{SGAresult}{命题}{2.10}
设 $X$ 是\Term{algebraic-space}{代数空间}，$U$ 是其\Term{open-subset}{开子集}，$Y$ 是 $U$ 在 $X$ 中的补集。则自然\Term{homomorphism}{同态}序列
\SourcePage{66}
\[
K(Y)\longrightarrow K(X)\longrightarrow K(U)\longrightarrow0\tag{2.30}
\]
\Term{exact}{正合}。
\end{SGAresult}
\textbf{注意。} 可与命题 2.3 比较。
\begin{proof}
$K(X)\to K(U)$ 为\Term{surjective-map}{满射}，立即由命题 2.5（或命题 2.6）得出。还须证明，若 $x\in K(X)$ 在 $K(U)$ 中的\Term{restriction}{限制}为零，则它在 $K(Y)$ 的像内。为此，回到 $K(X),K(U)$ 的显式定义：有 $X$ 上两个\Term{coherent-algebraic-sheaf}{凝聚代数层} $F,G$，使 $x=\gamma_X(F)-\gamma_X(G)$ 且 $\gamma_U(F)-\gamma_U(G)=0$。后一关系意味着，可找到 $U$ 上\Term{coherent-sheaf}{凝聚层} $P',Q'$，它们带有分商两两\Term{isomorphism}{同构}的\Term{filtration}{滤过}，并且在 $U$ 上有\Term{isomorphism}{同构}
\[
P=F+Q'\xrightarrow{\sim}Q=G+P'.
\]
由命题 2.5，\Term{sheaf}{层} $P',Q'$ 是 $X$ 上\Term{coherent-sheaf}{凝聚层}（仍以相同字母\Term{representation}{表示}）在 $U$ 上的\Term{restriction}{限制}；同样，由下面的\Term{lemma}{引理} 2.11，它们的\Term{filtration}{滤过}来自 $X$ 上相应\Term{sheaf}{层}的\Term{filtration}{滤过}。于是
\[
\gamma_X(F)-\gamma_X(G)=(\gamma_X(P)-\gamma_X(Q))+(\gamma_X(P')-\gamma_X(Q')).
\]
引入 $X$ 上 $P',Q'$ 的\Term{filtration}{滤过}的分商 $P'_i,Q'_j$，最后一项就是形如 $\gamma_X(P'_i)-\gamma_X(Q'_i)$ 的项之和，其中 $P'_i|_U\simeq Q'_i|_U$。因此归结为证明如下特例：若 $R,S$ 是 $X$ 上\Term{coherent-sheaf}{凝聚层}，且 $R|_U\simeq S|_U$，则 $\gamma_X(R)-\gamma_X(S)$ 在 $K(Y)$ 的像内。令 $T$ 为给定\Term{isomorphism}{同构} $R|_U\to S|_U$ 的图；它是 $(R\times S)|_U$ 的\Term{coherent-subsheaf}{凝聚子层}，故由\Term{lemma}{引理} 2.11，它是 $R\times S$ 的某个\Term{coherent-subsheaf}{凝聚子层}在 $U$ 上的\Term{restriction}{限制}，
\SourcePage{67}
仍将这个子\Term{sheaf}{层}记作 $T$。于是
\[
\gamma_X(R)-\gamma_X(S)=(\gamma_X(R)-\gamma_X(T))-(\gamma_X(S)-\gamma_X(T)),
\]
例如只须证明右端第一项在 $K(Y)$ 的像内。但\Term{projection}{投影} $u:T\to R$ 在 $U$ 上是\Term{isomorphism}{同构}，故
\[
\gamma_X(R)-\gamma_X(T)=\gamma_X(\operatorname{Coker}u)-\gamma_X(\operatorname{Ker}u),
\]
而 $\operatorname{Ker}u,\operatorname{Coker}u$ 的\Term{support}{支集}都在 $Y$ 内。
\end{proof}
现在只须证明下述\Term{lemma}{引理}。
\begin{SGAresult}{\Term{lemma}{引理}}{2.11}
设 $X$ 是\Term{algebraic-space}{代数空间}，$F$ 是 $X$ 上的\Term{coherent-algebraic-sheaf}{凝聚代数层}，$U$ 是 $X$ 的\Term{open-subset}{开子集}，$G$ 是 $F|_U$ 的\Term{coherent-subsheaf}{凝聚子层}。则 $G$ 是 $F$ 的某个\Term{coherent-subsheaf}{凝聚子层}在 $U$ 上的\Term{restriction}{限制}。
\end{SGAresult}
\begin{proof}
令 $\overline G$ 为 $F$ 的如下子\Term{sheaf}{层}：在 $X$ 的\Term{open-subset}{开集} $V$ 上，其\Term{section}{截面}是那些\Term{restriction}{限制}到 $U\cap V$ 后成为 $G$ 的\Term{section}{截面}的 $F$ 在 $V$ 上的\Term{section}{截面}。只须证明 $\overline G$ \Term{coherent}{凝聚}。可设 $X,U$ 都\Term{affine}{仿射}（因为 $U$ 是\Term{affine-open}{仿射开集} $U_i$ 的并，而 $\overline G$ 是相应各\Term{sheaf}{层} $\overline{G|_{U_i}}$ 的交）。于是 $U$ 的\Term{coordinate-ring}{坐标环}为 $X$ 的\Term{coordinate-ring}{坐标环} $A$ 的\Term{fraction-ring}{分式环} $A_S$；立即可验证，若 $F$ 由 $A$-模 $M$ 定义，则 $F|_U$ 由 $A_S$-模 $M_S=A_S\otimes_A M$ 定义。$F|_U$ 的子\Term{sheaf}{层} $G$ 便由 $M_S$ 的某个\Term{submodule}{子模} $N'$ 定义，而 $\overline G$ 恰由 $M$ 的\Term{submodule}{子模} $N$ 定义，$N$ 是 $N'$ 在\Term{canonical}{典范}\Term{homomorphism}{同态} $M\to M_S$ 下的\Term{inverse-image}{逆像}。
\end{proof}
\begin{SGAresult}{命题 2.10 的推论}{1}
设 $x\in K(X)$。$x$ 的\Term{filtration}{滤过} $\leq i$ 的充要条件是：存在 $X$ 中\Term{dimension}{维数} $\leq n-i$ 的\Term{closed-subset}{闭子集} $Y$，使 $x$ 在 $X-Y$ 上的\Term{restriction}{限制}为零。
\end{SGAresult}
\SourcePage{68}
\begin{SGAresult}{推论}{2}
设 $X$ \Term{connected}{连通}，\Term{dimension}{维数}为 $n$，在\Term{dimension}{维数} $n-1$ 处没有\Term{singular-point}{奇点}。若 $D,D'$ 是 $X$ 上两个\Term{linear-equivalence}{线性等价}的\Term{divisor}{除子}，则 $\gamma(D)\equiv\gamma(D')\pmod {K(X)^2}$。
\end{SGAresult}
\begin{proof}
由推论 1，可设 $X$ \Term{nonsingular}{非奇异}。此推论于是由\Term{generic}{一般}公式
\[
\gamma(D)\equiv1-\gamma(L(-D))\pmod {K(X)^2}\tag{2.31}
\]
得出，其中 $L(-D)$ \Term{representation}{表示}\Term{divisor}{除子} $-D$ 所定义的\Term{line-bundle}{秩 $1$ 向量丛}。为验证 (2.31)，将其写为
\[
\gamma(L(-D))\equiv1-\gamma(D)\pmod {K(X)^2}.\tag{2.31 bis}
\]
注意两端都是从\Term{divisor}{除子}群到环 $K(X)/K(X)^2$ 的\Term{invertible-element}{可逆元}\Term{multiplicative-group}{乘法群}的\Term{homomorphism}{同态}（注意 $K(X)^1/K(X)^2$ 的平方为零；使用推论 1，由显式公式 (2.15) 很容易验证这一点）。从而归结为 $D$ 是\Term{irreducible}{不可约}\Term{hypersurface}{超曲面}的情形。更\Term{generic}{一般}地，若 $D$ 是各\Term{branch}{分支}互\Term{disjoint}{不相交}的\Term{hypersurface}{超曲面}，则有等式
\[
\gamma(D)=1-\gamma(L(-D)),\tag{2.32}
\]
它来自\Term{sheaf}{层}的\Term{exact-sequence}{正合列}
\[
0\longrightarrow\mathcal O_X(L(-D))\longrightarrow\mathcal O_X\longrightarrow\mathcal O_D\longrightarrow0.
\]
\end{proof}
\begin{SGAresult}{\Term{theorem}{定理}}{2.12}
设 $X$ 是\Term{nonsingular}{非奇异}\Term{quasi-projective}{拟射影}簇，$Z,Z'$ 是 $X$ 上两个\Term{linear-equivalence}{线性等价}的\Term{codimension}{余维} $p$ 的\Term{homogeneous}{齐次}\Term{algebraic-cycle}{代数闭链}。则
\[
\gamma_X(Z)\equiv\gamma_X(Z')\pmod {K(X)^{p+1}}.\tag{2.33}
\]
\end{SGAresult}
\begin{proof}
由假设，可找到 $X\times k$ 上\Term{dimension}{维数}为 $n-p+1$ 的\Term{algebraic-cycle}{代数闭链} $Y$（$n$ 为 $X$ 的\Term{dimension}{维数}），
\SourcePage{69}
以及两个点 $a,a'\in k$，使
\[
Z=f_a^{-1}(Y),\qquad Z'=f_{a'}^{-1}(Y),
\]
其中置 $f_u(x)=(x,u)$。由公式 (2.23)，有
\[
\gamma_X(Z)\equiv f_a^!(\gamma_{X\times k}(Y))\pmod {K(X)^{p+1}},
\]
而 $\gamma_X(Z')$ 也有类似公式。于是只须证明 $f_a^!=f_{a'}^!$，这立即由命题 2.9 得出。
\end{proof}
\begin{SGAresult}{推论}{1}
$K(X)$ 的\Term{filtration}{滤过}与其\Term{ring-structure}{环结构}相容，即任意整数 $i,j\geq0$ 都有 $K(X)^iK(X)^j\subset K(X)^{i+j}$。此外，有自然的\Term{graded-ring}{分次环}满\Term{homomorphism}{同态}
\[
\varphi:A(X)\longrightarrow G(K(X)),\tag{2.34}
\]
其中右端\Term{representation}{表示} $K(X)$ 对\Term{filtration}{滤过} $(K(X)^i)_{i\geq0}$ 的\Term{associated-graded-ring}{相伴分次环}；此\Term{homomorphism}{同态}由条件
\[
\varphi(Z^{n-p})=\gamma(Z^{n-p})\pmod {K(X)^{p+1}}
\]
定义。
\end{SGAresult}
\begin{proof}
对整数 $i+j$ 归纳，直接由\Term{theorem}{定理} 2.12、命题 2.6 和 2.7，以及 Chow \Term{theorem}{定理} 2.1 得出。
\end{proof}
\begin{SGAresult}{推论}{2}
设 $X,Y$ 是两个\Term{nonsingular}{非奇异}\Term{quasi-projective}{拟射影}簇，$f:X\to Y$ 为\Term{morphism}{态射}。则\Term{homomorphism}{同态} $f^!:K(Y)\to K(X)$ 与\Term{filtration}{滤过}相容，而且下图交换：
\[
\begin{tikzcd}
G(K(Y))\arrow[r,"G(f^!)"]&G(K(X))\\
A(Y)\arrow[u,"\varphi_Y"]\arrow[r,"f^*"]&A(X)\arrow[u,"\varphi_X"'].
\end{tikzcd}
\]
\end{SGAresult}
\SourcePage{70}
\begin{proof}
利用 $f$ 的图，将 $f$ \Term{resolution}{分解}为 $X$ 到 $X\times Y$ 的\Term{immersion}{浸入}与 $X\times Y$ 到 $Y$ 的\Term{projection}{投影}的复合。对\Term{projection}{投影}，此推论平凡；对\Term{immersion}{浸入}，按推论 1 的证明进行，用命题 2.8 代替命题 2.7。
\end{proof}
\begin{SGAresult}{命题}{2.13（“\Term{cellular-decomposition}{胞腔分解}”）}
设 $P$ 是\Term{algebraic-space}{代数空间}（可以\Term{irreducible}{不可约}，也可以\Term{reducible}{可约}），带有递增\Term{closed-subset}{闭子集}族 $P^0\subset P^1\subset\cdots\subset P^n=P$，对每个 $0\leq i\leq n$，$P^i-P^{i-1}$ 的\Term{connected-component}{连通分支}都是\Term{isomorphism}{同构}于 $k^i$ 的\Term{algebraic-space}{代数空间} $V_{i,\alpha}$（置 $P^{-1}=\varnothing$）。则：

 a) 对任意\Term{algebraic-space}{代数空间} $X$，由 $X,P$ 上\Term{sheaf}{层}的\Term{tensor-product}{张量积}\Term{derive}{导出}的自然映射 $K(X)\otimes K(P)\to K(X\times P)$ 为\Term{surjective-map}{满射}。

 b) 各 $\gamma_P(V_{i,\alpha})$ 构成群 $K(P)$ 的一组\Term{generating-set}{生成元}。
\end{SGAresult}
\begin{proof}
对 $n$ 归纳，使用命题 2.9 的推论和命题 2.10。（注意：遗憾的是，我不知道各 $\gamma_P(V_{i,\alpha})$ 是否构成 $K(P)$ 的基，也不知道 a) 中的\Term{homomorphism}{同态}是否为\Term{bijective-map}{双射}。后文不需要知道这些。）
\end{proof}
\begin{SGAresult}{\Term{theorem}{定理}}{2.14}
设 $x_i$（$1\leq i\leq q$）、$E_j$（$1\leq j\leq r$）是字母，$n_j$（$1\leq j\leq r$）是正整数。把\Term{signature}{符号} $\lambda^k(x_i)$（$k\geq1$）及 $\lambda^k(E_j)$（$1\leq k\leq n_j$）视为\Term{indeterminate}{不定元}，设 $R$ 是这些\Term{indeterminate}{不定元}的整系数\Term{polynomial}{多项式}。假定对任意\Term{graded-ring}{分次环} $A$，在 $R$ 中以 $\widetilde A$ 的元素 $\xi_i,\eta_j$ 分别代入 $x_i,E_j$，
\SourcePage{71}
并满足\Term{restriction}{限制} $\varepsilon(\xi_i)=0$ 及 $\eta_j=[n_j;1+\eta_j^{(1)}+\cdots+\eta_j^{(n_j)}]$（故 $k>n_j$ 时 $\eta_j^{(k)}=0$），总有
\[
R\bigl((\lambda^k(\xi_i)),(\lambda^k(\eta_j))\bigr)\in\widetilde A^p.
\]
那么，设 $X$ 是\Term{irreducible}{不可约}\Term{nonsingular}{非奇异}\Term{quasi-projective}{拟射影}簇，无论 $x'_i\in K(X)^1$ 与 $X$ 上\Term{dimension}{维数}为 $n_j$ 的\Term{vector-bundle}{向量丛} $E'_j$ 如何选择（$1\leq i\leq q,\ 1\leq j\leq r$），均有
\[
R\bigl((\lambda^k(x'_i)),(\lambda^k(\gamma(E'_j)))\bigr)\in K(X)^p.\footnotemark
\]
\end{SGAresult}
\footnotetext{容易证明：若将 $R$ 用 $\gamma^k(\xi_i)$、$\gamma^k(E_i-\varepsilon(E_i))$ \Term{representation}{表示}，则上述假设意味着 $R$ 的所有项权重都 $\geq p$（其中 $\gamma^k$ 的权重为 $k$）。因而 2.14 是第 X 讲 1.3.2 的特例。}
\begin{proof}
由第 2 节，各 $x'_i$ 本身是形如 $\gamma(E)$ 的元素之差，其中 $E$ 为\Term{vector-bundle}{向量丛}。另一方面，将 $X$ 视为某\Term{projective-space}{射影空间} $P$ 的\Term{locally-closed}{局部闭}\Term{subvariety}{子簇}，由第 2 节和命题 2.5 可知：对任意 $F\in\mathscr T(X)$，$n$ 充分大时，\Term{sheaf}{层} $F(n)$ 由其\Term{section}{截面}生成。若 $F$ 由\Term{vector-bundle}{向量丛} $E$ 定义，这说明 $E(n)$ 是某 \Term{grassmannian}{Grassmann 簇} $G$ 上的\Term{canonical}{典范}\Term{vector-bundle}{向量丛}经某\Term{morphism}{态射} $X\to G$ 的\Term{inverse-image}{逆像}。考虑相应的\Term{morphism}{态射} $X\to P\times G$，可见 $E$ 是 $P\times G$ 上某\Term{vector-bundle}{向量丛}的\Term{inverse-image}{逆像}。将此用于所有涉及的\Term{vector-bundle}{向量丛}，并使用\Term{theorem}{定理} 2.12 的推论 2，就归结到 $X$ 是若干 \Term{grassmannian}{Grassmann 簇}的乘积的情形。这时，结论将由下述\Term{lemma}{引理}得出。
\end{proof}
\begin{SGAresult}{\Term{lemma}{引理}}{2.15}
设 $X$ 是若干 \Term{grassmannian}{Grassmann 簇}的乘积。则可找到从 $K(X)$ 到形如 $\widetilde A$ 的\Term{torsion-free}{无挠} \Term{lambda-ring}{$\lambda$-环}的 $\lambda$-\Term{homomorphism}{同态}（$A$ 为\Term{graded-ring}{分次环}），
\SourcePage{72}
它与\Term{filtration}{滤过}相容，并使相应\Term{homomorphism}{同态} $G(K(X))\to G(\widetilde A)$ 为\Term{injective-map}{单射}。
\end{SGAresult}
\begin{proof}
证明此\Term{lemma}{引理}将使用 Chow 的\Term{chern-class}{陈类}理论；若能直接证明本\Term{lemma}{引理}或\Term{theorem}{定理} 2.14，就可\Term{pushout}{推出}一种纯粹用\Term{sheaf}{层}建立的\Term{chern-class}{陈类}理论（见第 5 节）。

考虑\Term{homomorphism}{同态} $\widetilde C:K(X)\to\widetilde{A(X)}$，并置 $A=A(X)$；须证明\Term{homomorphism}{同态}
\[
\psi=G(\widetilde C):G(K(X))\longrightarrow G(\widetilde{A(X)})\simeq A(X)
\]
为\Term{injective-map}{单射}。$\widetilde C$ 与\Term{filtration}{滤过}相容，这对任意\Term{nonsingular}{非奇异}\Term{quasi-projective}{拟射影}簇都成立，并容易由命题 2.3 的推论得出。另一方面，$X$ 是若干 \Term{grassmannian}{Grassmann 簇} $G_j$ 的乘积，故 $X$ 上有\Term{vector-bundle}{向量丛} $E_j$，它们是各因子上\Term{canonical}{典范}\Term{vector-bundle}{向量丛}的\Term{inverse-image}{逆像}。我们使用如下已知事实：

 a) 在 $X$ 上，\Term{algebraic-cycle}{代数闭链}的\Term{linear-equivalence}{线性等价}等同于\Term{numerical-equivalence}{数值等价}，且 $A(X)$ 是\Term{finite-rank}{有限秩}群；

 b) 各 $C^i(E_j)$ 生成环 $A(X)$。

由 a) 及下述\Term{lemma}{引理}可知 $\varphi:A(X)\to G(K(X))$ 为\Term{injective-map}{单射}（故为\Term{bijective-map}{双射}）。
\end{proof}
\begin{SGAresult}{\Term{lemma}{引理}}{2.15}
设 $X^n$ 是\Term{nonsingular}{非奇异}\Term{projective-variety}{射影簇}，$Z^{n-p}$ 是 $X$ 上的\Term{algebraic-cycle}{代数闭链}，且 $\varphi(Z)\in G^p(K(X))$ 为零；则 $Z^{n-p}$ \Term{numerical-equivalence}{数值等价}于零。
\end{SGAresult}
\begin{proof}
只须使用\Term{additive-function}{加性函数} $\chi(F)=\sum_{i\geq0}(-1)^i\dim H^i(X,F)$，以及对每个 $x\in X$，此函数在 $F=\mathcal O_{(x)}$ 时等于 $1$ 这一事实。
\end{proof}
\SourcePage{73}
因此，在所考虑的情形，$A(X),G(K(X))$ 在每个\Term{degree}{次数}都有相同的\Term{finite-rank}{有限秩}，且 $G(K(X))$ \Term{torsion-free}{无挠}，因为 $A(X)$ 如此。由上面的 b)“\Term{form}{形式}地”得到：$\widetilde C:K(X)\to\widetilde{A(X)}$ 模\Term{finite-group}{有限群}为\Term{surjective-map}{满射}。由于 $G(K(X)),A(X)$ 在每个\Term{degree}{次数}有相同的\Term{finite-rank}{有限秩}，容易得出\Term{homomorphism}{同态} $G(\widetilde C):G(K(X))\to G(\widetilde{A(X)})$ 在每个\Term{degree}{次数}模\Term{finite-group}{有限群}为\Term{bijective-map}{双射}（对\Term{degree}{次数}归纳）。而 $G(K(X))$ \Term{torsion-free}{无挠}，所以\Term{homomorphism}{同态} $G(\widetilde C)$ 为\Term{injective-map}{单射}。证毕。

\textbf{注意。} \Term{homomorphism}{同态} $\widetilde C$ \Term{generic}{一般}并非\Term{surjective-map}{满射}，即使 $X$ 是二维\Term{projective-space}{射影空间}也如此。

\SGAsection{5}{陈类的层论定义及其在单射态射研究中的应用}
\begin{SGAresult}{\Term{theorem}{定理}}{2.16}
设 $X$ 是\Term{nonsingular}{非奇异}\Term{quasi-projective}{拟射影}簇。则对每个 $x\in K(X)$ 及每个整数 $i\geq0$，元素 $\gamma^i(x-\varepsilon(x))$ 的\Term{filtration}{滤过} $\geq i$。以 $c^i(x)\in G^i(K(X))$ \Term{representation}{表示}此元素模 $K(X)^{i+1}$ 的类；若置
\[
\widetilde c(x)=\left[\varepsilon(x),1+\sum_{i\geq0}c^i(x)\right]\in\widetilde{G(K(X))},
\]
则 $\widetilde c$ 满足与\Term{theorem}{定理} 2.2 中类似的条件。
\end{SGAresult}
\begin{proof}
第一断言，以及 $x\mapsto\widetilde c(x)$ 为 $\lambda$-\Term{ring-homomorphism}{环同态}这一事实，都由\Term{theorem}{定理} 2.14 得出（公式 (1.25) 蕴含 $\gamma_t(x+y)=\gamma_t(x)\gamma_t(y)$，这已表明 $\widetilde c$ 是加法\Term{homomorphism}{同态}）。$\widetilde c$ 的“\Term{functoriality}{函子性}”立即可验证；
\SourcePage{74}
(2.8) 由 (2.31) 得出。
\end{proof}
\begin{SGAresult}{推论}{1}有 $c^i(x)=\varphi(C^i(x))$。\end{SGAresult}
\begin{proof}
这来自\Term{chern-class}{陈类}的唯一性\Term{theorem}{定理}；如 Hirzebruch 的方法一样，通过过渡到\Term{flag-variety}{旗簇}证明。只须证明，若 $Y$ 是与 $X$ 上某\Term{vector-bundle}{向量丛}相伴的\Term{flag-variety}{旗簇}，则 $G(K(X))\to G(K(Y))$ 为\Term{injective-map}{单射}。这可由从 $K(X)$ 出发显式确定 $K(Y)$ 及其全部结构来验证（这里不给出细节）。
\end{proof}
\begin{SGAresult}{推论}{2}
假设对每个 $Z\in A(X)$，都存在从 $X$ 到某 \Term{grassmannian}{Grassmann 簇}乘积 $X'$ 的\Term{morphism}{态射} $f$，使 $Z$ 在 $f^*$ 的像内。则\Term{canonical}{典范}\Term{homomorphism}{同态}
\[
\varphi:A(X)\longrightarrow G(K(X)),\qquad\psi:G(K(X))\longrightarrow A(X)
\]
在\Term{dimension}{维数} $i$ 处满足关系
\[
\psi\varphi=(-1)^{i-1}(i-1)!\,1,\qquad
\varphi\psi=(-1)^{i-1}(i-1)!\,1.\tag{2.35}
\]
\end{SGAresult}
\begin{proof}
由于 $\varphi$ 为\Term{surjective-map}{满射}，显然只须证明第一式；为此可设 $X$ 是 \Term{grassmannian}{Grassmann 簇}的乘积。这时 $\varphi$ 为\Term{injective-map}{单射}，故只须证明 (2.35) 的第二式，由推论 1，可将它写为
\[
\gamma^i(x-\varepsilon(x))\equiv(-1)^{i-1}(i-1)!\,x\pmod {K(X)^{i+1}}
\]
对每个 $x\in K(X)^i$ 成立。$\psi$ 为\Term{injective-map}{单射}，故只须对 $\widetilde C(x)\in\widetilde{A(X)}^i$ 证明类似关系，而这就是 (1.22 bis) 的推论。
\end{proof}
\textbf{注记。} 推论 2 使人相信 (2.35) 在所有情形都成立；应当有一个直接证明。无论如何，
\SourcePage{75}
模去 $G(K(X))$ 的\Term{torsion-subgroup}{挠子群}后，很容易证明 (2.35) 的第二式；在\Term{characteristic}{特征} $0$ 时，同一公式又是 \Term{riemann-roch-theorem}{Riemann--Roch 定理}的推论。(2.35) 的意义（在它成立时）在于它说明 $\varphi$ 模\Term{torsion-group}{挠群}为\Term{injective-map}{单射}（\Term{torsion-group}{挠群}也许总是零！\footnote{并非如此，见原文第 30 页的注。}），因此将\Term{chern-class}{陈类}定义为 $G(K(X))$ 的元素，“损失并不大”；若像后文处理 \Term{riemann-roch-theorem}{Riemann--Roch 定理}那样与 $\mathbb Q$ 作\Term{tensor-product}{张量积}，则\Term{chern-class}{陈类}的两种定义便相同。

仍设 $X$ \Term{quasi-projective}{拟射影}且\Term{nonsingular}{非奇异}；再设 $Y$ 是 $X$ 的\Term{nonsingular}{非奇异}\Term{subvariety}{子簇}，$i:Y\to X$ 为包含\Term{injective-map}{单射}。对于 $i^!,i_!$，有下列公式：
\[
\begin{aligned}
\text{(i)}\quad &i^!\text{ 是 }\lambda\text{-环同态};\\
\text{(ii)}\quad &i_!(x\cdot i^!(y))=i_!(x)\cdot y,\\
&\text{特别地 }i_!i^!(y)=\xi y,\quad \xi=i_!i^!(1)=\gamma_X(Y);\\
\text{(iii)}\quad&i^!(i_!(x))=x\cdot\lambda_{-1}(N),\\
&\text{特别地 }i^!(\xi)=i^!i_!(1)=\lambda_{-1}(N);\\
\text{(iv)}\quad&i_!(x)i_!(x')=i_!(xx'\lambda_{-1}(N)).
\end{aligned}
\]
这里按定义，$N$ 是 $Y$ 在 $X$ 中的\Term{conormal-bundle}{余法丛}在 $K(Y)$ 中的类。公式 (i)、(ii) 对任意\Term{morphism}{态射}都成立，无需是\Term{injective-map}{单射}；(iii)、(iv) 则由定义 (2.15) 和\Term{inverse-image}{逆像}的类似公式验证（取 $x,x'$ 为\Term{vector-bundle}{向量丛}的类，并使用事实
\[
\operatorname{Tor}^{\mathcal O_X}_i(\mathcal O_Y,\mathcal O_Y)=\mathcal O_Y(\lambda^iN).\tag{2.36}
\]
）
\SourcePage{76}
公式 (iv) 表明，若在 $K(Y)$ 上定义新的\Term{multiplication}{乘法}结构
\[
x\mathbin{\cdot_N}x'=xx'\lambda_{-1}(N)
\]
（参见第一章第 4 节），则 $i_!$ 为\Term{ring-homomorphism}{环同态}。还有补充公式
\[
\text{(v)}\qquad i_!(\lambda^p(N,x))=\lambda^p(i_!(x))\qquad(p\geq1),
\]
意思是 $i_!$ 定义从 $K(Y)_N$ 到 $K(X)$ 的 $\lambda$-\Term{homomorphism}{同态}。公式 (v) 在所有情形下无疑都成立，但目前只在\Term{characteristic}{特征} $0$ 时得到证明，因为它是\Term{theorem}{定理} 1.4 和 2.8 的联合推论。

对 $i_*,i^*$ 有类似公式：
\[
\begin{aligned}
\text{(i bis)}\quad&i^*\text{ 是分次环同态};\\
\text{(ii bis)}\quad&i_*(x\cdot i^*(y))=i_*(x)\cdot y,\\
&\text{特别地 }i_*i^*(y)=\xi'y,\quad\xi'=i_*(1);\\
\text{(iii bis)}\quad&i^*(i_*(x))=x\cdot(-1)^qN^q,\\
&\text{特别地 }i^*(\xi')=(-1)^qN^q;\\
\text{(iv bis)}\quad&i_*(x)i_*(x')=i_*(xx'\cdot(-1)^qN^q);\\
\text{(v bis)}\quad&i_*\text{ 使次数增加 }q.
\end{aligned}
\]
在这些公式中，$q$ 是 $Y$ 在 $X$ 中的\Term{codimension}{余维}，并置 $N^i=C^i(N)$。只有 (iii bis)、(iv bis) 需要证明；我不知道它们是否已广为人知（甚至不知道它们是否正确！！！）。这里只指出，若在 $G(K(Y)),G(K(X))$ 中而非 $A(Y),A(X)$ 中考虑，则它们由公式 (ii)、(iv) 取商得到。这里使用 $i_*,i^*$ 分别由 $i_!,i^!$ 取商而来，并注意 $i_!$ 使\Term{filtration}{滤过}移动 $q$。
\SourcePage{77}
\textbf{注意。} 连接 $i_!,i^!$（分别为 $i_*,i^*$）的一切“\Term{universal}{普遍}”公式都是上述公式的\Term{form}{形式}推论；这不难看出。

由公式 (v) 得到
\[
i_!(\gamma^p(N,x))=\gamma^p(i_!(x)).
\]
\Term{theorem}{定理} 2.14 和命题 1.5 联合表明：当 $p=q+j$ 时，$\gamma^p(N,x)$ 的\Term{filtration}{滤过} $\geq j$；若以 $G_p(N,x)$ \Term{representation}{表示}其在 $G^j(K(Y))$ 中的类，则\Term{theorem}{定理} 2.16 给出
\[
i_*(G_p(N,x))=c^p(i_!(x)).
\]
最后，由命题 1.5 和\Term{theorem}{定理} 2.14，可将 $G_p(N,x)$ 表为各 $c^i(N),c^i(x),\varepsilon(x)$ 的整系数\Term{polynomial}{多项式}。由此得到：
\begin{SGAresult}{\Term{theorem}{定理}}{2.17}
设域 $k$ 的\Term{characteristic}{特征}为 $0$，$i:Y\to X$ 是\Term{nonsingular}{非奇异}\Term{quasi-projective}{拟射影}簇之间的\Term{injective-map}{单射}\Term{morphism}{态射}。则公式 (i) 至 (v) 成立；此外，对每个整数 $j\geq0$，有公式
\[
c^{q+j}(i_!(x))=i_*\left(\frac{\widetilde c(x)*\lambda_{-1}(\widetilde c(N))}{(-1)^qN^q}\right)^{(j)}.\tag{2.37}
\]
\end{SGAresult}
公式 (2.37) 右端是一种简写，\Term{representation}{表示}命题 1.5 所述的关于 $\widetilde c(x),\widetilde c(N)$ 的\Term{set-component}{分量}的\Term{polynomial}{多项式}。
\begin{SGAresult}{推论}{}
在上述条件下，有
\[
\operatorname{ch}(\widetilde c(i_!(x)))=i_*\bigl(\operatorname{ch}(\widetilde c(x))\mathscr T(\check N)^{-1}\bigr),\tag{2.38}
\]
其中 $\operatorname{ch}$ 是公式 (1.26) 定义的\Term{chern-homomorphism}{陈同态}。
\end{SGAresult}
\SourcePage{78}
\begin{proof}
只须将 (2.37) 写为
\[
\widetilde c(i_!(x))=1+i_*\left(\frac{\widetilde c(x)*\lambda_{-1}(\widetilde c(N))-1}{(-1)^qN^q}\right),\tag{2.37 bis}
\]
再使用显式公式 (1.29)、公式 (iv bis)，最后使用 (1.30)。
\end{proof}
可以把 (2.38) 写为
\[
\operatorname{ch}_X(F)=i_*\bigl(\operatorname{ch}_Y(F)\mathscr T(T_{X/Y})^{-1}\bigr).\tag{2.38 bis}
\]
其中 $F$ 是 $Y$ 上的\Term{coherent-algebraic-sheaf}{凝聚代数层}（左端将它视为 $X$ 上在 $Y$ 外为零的\Term{sheaf}{层}），$\operatorname{ch}_X,\operatorname{ch}_Y$ 分别是在 $X,Y$ 上取的\Term{chern-homomorphism}{陈同态}，$T_{X/Y}$ 是 $Y$ 在 $X$ 中的\Term{normal-bundle}{法丛}。

\textbf{注记。} 与 (2.38) 不同，(2.37) 是一个“无分母公式”。即使在\Term{characteristic}{特征} $0$ 时，目前也只在 $G(K(X))$ 中证明了它，在 $A(X)$ 中尚未证明，两者也许并不相同。它无疑在任意\Term{characteristic}{特征}及 $A(X)$ 中都成立。注意，立即可验证两端经 $i^*$ 的像相同；再\Term{action}{作用} $i_*$，得到它们的差与 $N^q$ 的乘积为零。这提示一种完全不同的证明方法：尝试将 $Y$ 实现为某个簇 $X'$ 中\Term{subvariety}{子簇} $Y'$ 的\Term{transverse}{横截}\Term{inverse-image}{逆像}，并要求在\Term{dimension}{维数} $\leq n=\dim X$ 时，$Y'$ 在 $A(X')$ 中的类不是\Term{zero-divisor}{零因子}……遗憾的是，不能期望 $X'$ 为 \Term{grassmannian}{Grassmann 簇}或与之类似的\Term{object}{对象}。

注意，若忽略\Term{torsion-element}{挠元素}，公式 (2.38) 与 (2.37) \Term{equivalence}{等价}。

\SourcePage{79}
\SGAsection{6}{Riemann--Roch 定理}
其陈述如下：设 $f:Y\to X$ 是\Term{nonsingular}{非奇异}\Term{quasi-projective}{拟射影}簇到同类簇的\Term{proper}{固有}\Term{morphism}{态射}。则
\[
f_*(\operatorname{ch}_Y(x)\mathscr T(Y))=\operatorname{ch}_X(f_!(x))\mathscr T(X).\tag{2.39}
\]
（两端均属于 $A(X)\otimes\mathbb Q$，且 $\mathscr T(X)=\mathscr T(C(T_X))$。）

此\Term{theorem}{定理}允许在已知 $x,T_X,T_Y$ 的\Term{chern-class}{陈类}及\Term{homomorphism}{同态} $f_*$ 时，计算 $f_!(x)$ 的\Term{chern-class}{陈类}，精确到\Term{torsion-element}{挠元素}。或者说，若借助 $\operatorname{ch}_Y,\operatorname{ch}_X$，把 $K(Y)\otimes\mathbb Q,K(X)\otimes\mathbb Q$ 分别与 $A(Y)\otimes\mathbb Q,A(X)\otimes\mathbb Q$ 等同（如果 (2.35) 成立的话！），则可利用 $f_*$ 及 $T_X,T_Y$ 的\Term{chern-class}{陈类}确定\Term{homomorphism}{同态} $f_!$。

我只知道如何在\Term{characteristic}{特征} $0$ 时证明 (2.39)\footnote{下文有一个适用于任意\Term{characteristic}{特征}、依据不同原理的证明（原文第 51 页）。}，且证明所得是在 $G(K(X))\otimes\mathbb Q$ 中的等式。为此，将 $Y$ 视为\Term{projective-space}{射影空间} $P$ 的\Term{locally-closed}{局部闭}\Term{subvariety}{子簇}，把 $f$ \Term{resolution}{分解}为 $Y$ 到 $P\times X$ 的\Term{injective-map}{单射}\Term{morphism}{态射}及 $P\times X$ 到 $X$ 的\Term{projection}{投影}\Term{morphism}{态射}的复合，并分别为这两个\Term{morphism}{态射}证明 (2.39)。

对于\Term{injective-map}{单射}\Term{morphism}{态射}，立即可见 (2.38) 和 (2.39) \Term{equivalence}{等价}（\Term{characteristic}{特征}为 $0$ 的假设只通过\Term{theorem}{定理} 2.8 用到；该\Term{theorem}{定理}使我们能用\Term{coherent-algebraic-sheaf}{凝聚代数层} $F$ 的\Term{local-scheme}{局部}\Term{cohomology}{上同调}\Term{invariants}{不变量}\Term{representation}{表示} $\lambda^i(\gamma(F))$）。

对于\Term{projection}{投影}\Term{morphism}{态射} $f:P\times X\to X$，由命题 2.13，可设 $x$ 形如 $y\otimes z$，其中 $y\in K(P)$，
\SourcePage{80}
$z\in K(X)$。于是 \Term{kunneth-formula}{Künneth 公式}给出 $f_!(x)=\chi(y)z$，这里，对 $P$ 上\Term{coherent-algebraic-sheaf}{凝聚代数层} $F$，置
\[
\chi(F)=\sum_{i\geq0}(-1)^i\dim H^i(P,F).
\]
立即归结到证明公式
\[
\chi(F)=\kappa_r(\operatorname{ch}_P(F)\mathscr T(P)),\tag{2.40}
\]
其中 $\kappa_r$ \Term{representation}{表示}\Term{degree}{次数} $r=\dim P$ 的\Term{set-component}{分量}。由命题 2.13，可设 $F=\mathcal O_Q$，其中 $Q$ 为 $P$ 的\Term{linear}{线性}\Term{subvariety}{子簇}。于是 $\chi(F)=1$，因为这是 $Q$ 的\Term{arithmetic-genus}{算术亏格}；只须说明 (2.40) 右端等于 $1$。

事实上，若 $V$ 是 $P$ 的\Term{hyperplane}{超平面}，且置 $\xi=\gamma_P(\mathcal O_V)$、$L=\ell(V)$，则有
\[
1-\xi=\gamma_P(E(-L))
\]
（其中 $E(-L)$ 是由 $-L$ 定义的\Term{vector-bundle}{向量丛}），故
\[
C_P(1-\xi)=[1,1-L],
\]
从而
\[
\operatorname{ch}_P(1-\xi)=e^{-L},\qquad
\operatorname{ch}_P(\xi)=1-e^{-L},
\]
所以
\[
\operatorname{ch}_P(\xi^p)=(1-e^{-L})^p.
\]
若 $p$ 是 $Q$ 在 $P$ 中的\Term{dimension}{维数}，则 $\xi^p=\gamma_P(Q)$，因此
\[
\operatorname{ch}_P(F)\mathscr T(P)=\left(\frac{L}{1-e^{-L}}\right)^{r+1}(1-e^{-L})^p.
\]
必须验证此关于 $L$ 的\Term{formal-power-series}{形式幂级数}中 $L^r$ 的系数为 $1$。这来自\Term{formal-power-series}{形式幂级数} $L/(1-e^{-L})$ 的\Term{characteristic}{特征}性质，即 $(L/(1-e^{-L}))^{q+1}$ 中 $L^q$ 的系数为 $1$（令 $q=r+1-p$）。

\SourcePage{81}
\textbf{参考文献}
\begin{enumerate}
\item[{[1]}] Hirzebruch：《\Term{algebraic-geometry}{代数几何}中的新\Term{topology}{拓扑}方法》（\textit{Neue topologischen Methoden in der algebraischen Geometrie}）。
\item[{[2]}] Serre：《\Term{coherent-algebraic-sheaf}{凝聚代数层}》（\textit{Faisceaux algébriques cohérents}），\textit{Ann. of Maths.}。
\item[{[3]}] Grothendieck：《论\Term{coherent-algebraic-sheaf}{凝聚代数层}》（\textit{Sur les faisceaux algébriques cohérents}），载 Cartan 讨论班，1956/57。
\end{enumerate}

\SGAsection{}{Riemann--Roch 定理的证明}
\begin{center}（Grothendieck，1957 年 11 月 1 日）\end{center}
\textbf{* 记号回顾}\footnote{星号 *\ $\cdots$\ * 之间的段落为抄录者插入的文字。}

[采用《Riemann--Roch 报告》中的记号，下文将该报告引用为 R.R.R.。]

所考虑的簇一律\Term{nonsingular}{非奇异}且\Term{quasi-projective}{拟射影}（即\Term{isomorphism}{同构}于某\Term{projective-space}{射影空间}的\Term{locally-closed}{局部闭}\Term{subvariety}{子簇}）。\Term{base-field}{基域}\Term{algebraically-closed}{代数闭}，\Term{characteristic}{特征}任意。

对这样的簇 $X$，以 $K(X)$ \Term{representation}{表示}由 $X$ 上\Term{coherent-sheaf}{凝聚层}的类生成的\Term{abelian-group}{阿贝尔群}（模去将\Term{extension}{扩张}与\Term{direct-sum}{直和}等同的关系）；Tor 的\Term{alternating-sum}{交错和}使 $K(X)$ 成为环。若 $F$ 是\Term{sheaf}{层}（分别地，$E$ 是\Term{vector-bundle}{向量丛}，$Y$ 是 $X$ 的\Term{subvariety}{子簇}），
\SourcePage{82}
以 $\gamma_X(F)$（分别为 $\gamma_X(E),\gamma_X(Y)$）\Term{representation}{表示} $K(X)$ 中对应于 $F$（分别为 $E,\mathcal O_Y$）的元素。（实际上常直接写 $F,E$ 来代替 $\gamma_X(F),\ldots$。）

以 $A(X)$ \Term{representation}{表示} $X$ 的\Term{algebraic-cycle}{代数闭链}类关于\Term{rational-equivalence}{有理等价}的\Term{graded-ring}{分次环}（参见 Chow）；$A^i(X)$ 由\Term{dimension}{维数}为 $i$ 的\Term{algebraic-cycle}{代数闭链}组成。

若 $E$ 是\Term{vector-bundle}{向量丛}，以 $C(E)=\sum C^i(E)$，$C^i(E)\in A^i(X)$，\Term{representation}{表示}其\Term{chern-class}{陈类}（定义为 Schubert \Term{algebraic-cycle}{代数闭链}的\Term{inverse-image}{逆像}）；由\Term{linear}{线性}性得出每个 $y\in K(X)$ 的 $C^i(y)$。若 $C(y)=\prod(1+a_i)$，依 Hirzebruch 置
\[
\operatorname{ch}(y)=\sum e^{a_i},\qquad T(y)=\prod\frac{a_i}{1-e^{-a_i}}.
\]
它们是 $A(X)\otimes\mathbb Q$ 的元素。当 $y$ 是 $X$ 的\Term{tangent-bundle}{切丛}时，写 $T(X)$ 代替 $T(y)$。

若 $G$ 是\Term{vector-bundle}{向量丛}，以 $\lambda^pG$ \Term{representation}{表示}其第 $p$ \Term{exterior-power}{外幂}，并置 $\lambda_{-1}(G)=\sum(-1)^p\lambda^p(G)$；这是 $K(X)$ 的元素。

若 $f:Y\to X$ 是簇之间的\Term{proper}{固有}\Term{regular-map}{正则映射}，定义加法\Term{homomorphism}{同态}
\[
f^*:A(X)\longrightarrow A(Y),\qquad f_*:A(Y)\longrightarrow A(X).
\]
前者\Term{homogeneous}{齐次}，\Term{degree}{次数}为 $0$，且为\Term{multiplication}{乘法}\Term{homomorphism}{同态}；后者\Term{homogeneous}{齐次}，\Term{degree}{次数}为 $\dim X-\dim Y$。

还定义\Term{homomorphism}{同态}
\[
f^!:K(X)\longrightarrow K(Y),\qquad f_!:K(Y)\longrightarrow K(X).
\]
若 $F$ 是 $X$ 上的\Term{coherent-sheaf}{凝聚层}，按定义，$f^!(F)$ 等于各 $\operatorname{Tor}^{\mathcal O_X}_p(\mathcal O_Y,F)$ 的\Term{alternating-sum}{交错和}。对于\Term{vector-bundle}{向量丛}，这就是\Term{inverse-image}{逆像}运算。
\SourcePage{83}
若 $G$ 是 $Y$ 上\Term{coherent-sheaf}{凝聚层}，$f_*(G)$ 在 $X$ 上\Term{coherent}{凝聚}。以 $R^pf_*$ \Term{representation}{表示}如此定义的\Term{functor}{函子} $f_*$ 的第 $p$ \Term{derived-functor}{导出函子}，则按定义
\[
f_!(G)=\sum(-1)^pR^pf_*(G).\qquad *
\]
\begin{SGAresult}{\Term{riemann-roch-theorem}{Riemann--Roch 定理}}{}
设 $f:Y\to X$ 是从\Term{nonsingular}{非奇异}\Term{quasi-projective}{拟射影}簇 $Y$ 到另一这样的簇 $X$ 的\Term{proper}{固有}\Term{morphism}{态射}（* \Term{regular-map}{正则映射} *），且 $y\in K(Y)$。则在 $A(X)\otimes\mathbb Q$ 中有
\[
f_*(\operatorname{ch}(y)T(Y))=\operatorname{ch}(f_!(y))T(X).\tag{1}
\]
\end{SGAresult}
\textbf{证明（* 简略 *）。}
\begin{SGAresult}{\Term{lemma}{引理}}{1}
考虑簇的\Term{proper}{固有}\Term{morphism}{态射} $Z\xrightarrow{g}Y\xrightarrow{f}X$。若 R.R. \Term{theorem}{定理}对 $f,g$ 成立，则对 $f\circ g$ 成立；若对 $fg,g$ 成立，则对 $f$ 成立。若 $f_*$ 为\Term{injective-map}{单射}，且 R.R. \Term{theorem}{定理}对 $fg,f$ 成立，则对 $g$ 成立。
\end{SGAresult}
\begin{proof}
平凡。（* 实际上，当假设 R.R. 对 $fg,g$ 成立时，还须假定 $g_!$ 为\Term{surjective-map}{满射}，才能断言 R.R. 对 $f$ 成立。*）
\end{proof}

\Term{proper}{固有}\Term{morphism}{态射}是一个\Term{injective-map}{单射}与一个\Term{projection}{投影} $X\times P\to X$ 的复合（$P$ 为\Term{projective-space}{射影空间}），所以归结为分别处理这两种情形。第二种情形，正如已经看到的，借助 R.R.R. 的命题 2.13（“\Term{cellular-decomposition}{胞腔分解}”），可归结为标准计算。现在设 $Y\subset X$，包含\Term{injective-map}{单射}\Term{morphism}{态射}为 $i$。公式 (1) 变成
\[
\operatorname{ch}(i_!(y))=i_*(\operatorname{ch}(y)T(E)^{-1}),\tag{2}
\]
其中 $E=T_{X/Y}$ 为 $Y$ 在 $X$ 中的\Term{normal-bundle}{法丛}。
\SourcePage{84}
\begin{SGAresult}{\Term{lemma}{引理}}{2}
假设 $X$ 上存在\Term{rank}{秩} $p=\operatorname{codim}_X Y$ 的\Term{vector-bundle}{向量丛} $G$，使
\[
\gamma_X(Y)=i_!(1)=\lambda_{-1}(G),\qquad i^!(G)=\check E,\qquad(-1)^pC^p(G)=Y.\tag{3}
\]
若
\[
y=i^!(x),\qquad x\in K(X),\tag{4}
\]
则公式 (2) 成立。
\end{SGAresult}
\begin{proof}
立即由 $\operatorname{ch}$ 为\Term{ring-homomorphism}{环同态}及公式
\[
\operatorname{ch}(\lambda_{-1}(G))=(-1)^pC^p(G)T(\check G)^{-1}\tag{5}
\]
得出。
\end{proof}
\begin{SGAresult}{推论}{}
若 $Y$ 为\Term{hypersurface}{超曲面}，且 $y$ 形如 $i^!(x)$，其中 $x\in K(X)$，则 R.R. 成立。
\end{SGAresult}
\begin{SGAresult}{\Term{lemma}{引理}}{3}
设 $X=Y\times P$，其中 $P$ 为\Term{projective-space}{射影空间}，且 $i$ 为包含\Term{injective-map}{单射} $u\mapsto(u,v_0)$，$v_0\in P$。则 $i_*$ 为\Term{injective-map}{单射}，且 R.R. 成立。
\end{SGAresult}
\begin{proof}这纯粹是标准计算。\end{proof}

再设 $X,Y,y$ 任意（仍有 $Y\subset X$）。设 $X'$ 是 $X$ 沿 $Y$ 的\Term{blowup}{爆破}所得的簇，$Y'$ 是 $Y$ 的\Term{inverse-image}{逆像}\Term{hypersurface}{超曲面}。有\Term{canonical-map}{典范映射}组成的图（$j$ 为\Term{canonical-map}{典范映射}）：
\[
\begin{tikzcd}
Y\arrow[r,"i"]&X\\
Y'\arrow[u,"g"]\arrow[r,"j"]&X'\arrow[u,"f"'].
\end{tikzcd}
\]
$E$ 在 $Y'$ 上的\Term{inverse-image}{逆像}有一个\Term{rank}{秩} $1$ \Term{subbundle}{子丛}，将其\Term{dual}{对偶}记为 $L$。此外，$L^{-1}$ 也是\Term{divisor}{除子} $Y'$ 在 $X'$ 上定义的丛 $L(Y')$ 在 $Y'$ 上的\Term{restriction}{限制}。置 $F=E/L^{-1}$。

初等计算给出下述\Term{lemma}{引理}。
\SourcePage{85}
\begin{SGAresult}{\Term{lemma}{引理}}{4}
有以下公式（其中 $p=\operatorname{codim}_X Y$）：
\[
f^!i_!(y)=j_!(\lambda_{-1}(\check F)y).\tag{6}
\]
\[
\lambda_{-1}(\check F)\equiv\gamma^{p-1}(E)\lambda^p(\check E)\pmod{(1-L)}.\tag{7}
\]\footnotemark
\[
f_*(1)=1,\qquad\text{因而 }f_*f^*=1.\tag{8}
\]
\[
f^*i_*(\alpha)\equiv j_*(\alpha C^{p-1}(F))\pmod{\operatorname{Ker}f_*}\quad\text{对每个 }\alpha\in A(Y).\tag{9}
\]
\end{SGAresult}
\footnotetext{这个公式无谓地复杂，例如可用第 VII 讲 4.1 代替。}
[为使书写简便，通过\Term{homomorphism}{同态} $g^!,g^*$ 令 $K(Y)$ \Term{action}{作用}在 $K(Y')$ 上、$A(Y)$ \Term{action}{作用}在 $A(Y')$ 上。* 公式 (7) 中的 $\gamma^{p-1}(E)$ 似乎\Term{representation}{表示} $\lambda^{p-1}(E-2)$，即 $\sum_{i+j=p-1}\binom{-p}{j}\lambda^i(E)$；同余发生在 $K(Y')$ 中。*]

公式 (6) 来自等式
\[
\operatorname{Tor}^{\mathcal O_X(v)}_i(\mathcal O_Y(v),\mathcal O_{Y'}(v'))=\lambda^i\check F(v')\qquad(v'\in Y',\ v=g(v')).\tag{6 bis}
\]
公式 (7) 来自以下公式：
\[
g_!(L^{-i})=0\quad(0<i\leq p-1),\qquad g_!(1)=1,\tag{10}
\]
这些公式本身是 \Term{kunneth-formula}{Künneth 公式}和\Term{projective-space}{射影空间}\Term{cohomology}{上同调}的推论。

公式 (8) 平凡；最后，(9) 来自 $g_*(C^{p-1}(F))=1$。后者例如可由如下显然的公式得出：
\[
g_*(\xi^i)=0\quad(0\leq i<p-1),\qquad g_*(\xi^{p-1})=1,\tag{10 bis}
\]
其中 $\xi=C^1(L)\in A^1(Y')$。

[实际上，(9) 极有可能是恒等式，而不只是同余式。]
\SourcePage{86}
由公式 (7) 立即得到 $y\lambda_{-1}(\check F)\in j^!(K(X'))$，只要
\[
y\gamma^{p-1}(E)\lambda^p(\check E)\in i^!(K(X)).\tag{11}
\]
假定条件 (11) 满足。由\Term{lemma}{引理} 2 的推论，R.R. 对 $(Y',X',y\lambda_{-1}(\check F))$ 成立。另一方面，由公式 (8)，为证明 (2)，只须证明
\[
f^*\operatorname{ch}(i_!(y))\equiv f^*\bigl(i_*(\operatorname{ch}(y)T(E)^{-1})\bigr)\pmod{\operatorname{Ker}f_*}.\tag{12}
\]
而这最后一个公式，使用 (6)、(9)、对 $(Y',X',y\lambda_{-1}(\check F))$ 的 R.R.，以及公式
\[
\operatorname{ch}(\lambda_{-1}(\check F))=C^{p-1}(F)T(F)^{-1}\qquad\text{（参见公式 (5)）},\tag{13}
\]
立即可验证。

这样，只要 $y\gamma^{p-1}(E)=0$，R.R. 就得到证明；特别地（因为 $\gamma^{p-1}(E)$ 的\Term{filtration}{滤过} $\geq p-1$，* 这里在 $K(Y)$ 上按\Term{support}{支集}的\Term{codimension}{余维}取\Term{filtration}{滤过} *），只要 $p-1>\dim Y$，即
\[
\dim Y<\frac{\dim X-1}{2},
\]
便成立。结合\Term{lemma}{引理} 1 和\Term{lemma}{引理} 3，可归结到这种情形（取\Term{projective-space}{射影空间} $P$ 的\Term{dimension}{维数} $r$ 充分大，例如 $r>\dim X+1$）。证毕。

\textbf{注意。} 上面用了 $\gamma^{p-1}(E)$ 的\Term{filtration}{滤过} $\geq p-1$ 这一事实。可通过研究与 $E$ 相伴的\Term{projective-space}{射影空间}丛 $Y'$，直接证明这个事实，无须像 R.R.R. 2.16 那样经过 \Term{grassmannian}{Grassmann 簇}。更\Term{generic}{一般}地，这个方法给出 R.R.R. 的\Term{theorem}{定理} 2.14。

\SGAsection{}{Grothendieck 书信节录}
……“从实质上说”，新方法依赖于在下列两种情形确定 $K(X)$ 和 $A(X)$：$X$ 或者是与\Term{vector-bundle}{向量丛}相伴的\Term{projective-bundle}{射影丛}，或者是沿一个\Term{nonsingular}{非奇异}\Term{subvariety}{子簇}\Term{blowup}{爆破}而得的簇。
\SourcePage{87}
在第二种情形，我还没有完全确定，尚缺一点点。采用那篇短文的记号，须证明公式 (9)（但要作为等式）：
\[
f^*(i_*(\alpha))=j_*(\alpha C^{p-1}(F)),\qquad\alpha\in A(Y).\tag{9}
\]
可见两端经 $f_*$ 的像相同，经 $j^*$ 的像也相同，所以问题\Term{equivalence}{等价}于
\[
\operatorname{Ker}f_*\cap\operatorname{Ker}j^*=0.\tag{9 bis}
\]
无论如何，\Term{reduced}{约化}到 $G(K(X'))$ 后（这至多消去一些\Term{torsion-element}{挠元素}），(9) 成立，因为它此时由 (6) 得出；因此几乎可以肯定它成立。

……我还要告诉你：等式 (9) 很容易从 (6) 得出，无须经过麻烦的 (6 bis)……

\begin{flushright}核副本无误：J.-P. S.\end{flushright}
